% Source: https://artofproblemsolving.com/community/c3h1988121_a_proportion_problem
% Stable identity: aops-topic-1988121-q01
As fractions of volume, pump $a$ fuel is $\dfrac{1}{6}$ water, $\dfrac{1}{3}$ alcohol and $\dfrac{1}{2}$ gasoline, and pump $b$ fuel is $\dfrac{1}{4}$ water, $\dfrac{1}{3}$ alcohol and $\dfrac{5}{12}$ gasoline. If a fraction $\lambda$ ($0<\lambda<1$) of the mixture comes from pump $a$, the mixture is
\[
\text{alcohol } \dfrac{1}{3},\qquad \text{water } \dfrac{\lambda}{6}+\dfrac{1-\lambda}{4},\ \text{which lies strictly between } \dfrac{1}{6} \text{ and } \dfrac{1}{4} .
\]
Every option has alcohol $\dfrac{1}{3}$ of the total ($\dfrac{5}{15},\dfrac{5}{15},\dfrac{5}{15},\dfrac{6}{18},\dfrac{9}{27}$), so the water fraction decides:
\[
\text{a) }\dfrac{2}{15}<\dfrac{1}{6},\quad \text{b) }\dfrac{3}{15}=\dfrac{1}{5},\quad \text{c) }\dfrac{4}{15}>\dfrac{1}{4},\quad \text{d) }\dfrac{5}{18}>\dfrac{1}{4},\quad \text{e) }\dfrac{7}{27}>\dfrac{1}{4} .
\]
Only b) is possible. Indeed $\dfrac{\lambda}{6}+\dfrac{1-\lambda}{4}=\dfrac{1}{5}$ gives $\lambda=\dfrac{3}{5}$, and then the gasoline fraction is $\dfrac{3}{5}\cdot\dfrac{1}{2}+\dfrac{2}{5}\cdot\dfrac{5}{12}=\dfrac{7}{15}$. Mixing $3$ parts of fuel $a$ with $2$ parts of fuel $b$ gives $\boxed{\text{b) } 3:5:7}$.
